\begin{register}{H}{UART\_FIFO\_REG}{0x{}0000}\label{regdesc:UARTFIFOREG}
\regfield{(reserved)}{24}{8}{000000000000000000000000}%
\regfield{UART\_RXFIFO\_RD\_BYTE}{8}{0}{{0}}%
\reglabel{Reset}\regnewline%
\begin{regdesc}\begin{reglist}
\label{fielddesc:UARTRXFIFORDBYTE}\item [UART\_RXFIFO\_RD\_BYTE] UART\regindex{n} 通过此字段访问 FIFO。 (RO)
\end{reglist}\end{regdesc}
\end{register}


\begin{register}{H}{UART\_MEM\_CONF\_REG}{0x{}0060}\label{regdesc:UARTMEMCONFREG}
\regfield{(reserved)}{3}{29}{000}%
\regfield{UART\_MEM\_FORCE\_PU}{1}{28}{{0}}%
\regfield{UART\_MEM\_FORCE\_PD}{1}{27}{{0}}%
\regfield{UART\_RX\_TOUT\_THRHD}{10}{17}{{0x{}a}}%
\regfield{UART\_RX\_FLOW\_THRHD}{10}{7}{{0x{}0}}%
\regfield{UART\_TX\_SIZE}{3}{4}{{0x{}1}}%
\regfield{UART\_RX\_SIZE}{3}{1}{{1}}%
\regfield{(reserved)}{1}{0}{0}%
\reglabel{Reset}\regnewline%
\begin{regdesc}\begin{reglist}
\label{fielddesc:UARTRXSIZE}\item [UART\_RX\_SIZE] 配置 RAM 分配给 RX FIFO 的空间大小。 默认为 128 字节。 (R/W)
\label{fielddesc:UARTTXSIZE}\item [UART\_TX\_SIZE] 配置 RAM 分配给 TX FIFO 的空间大小。 默认为 128 字节。 (R/W)
\label{fielddesc:UARTRXFLOWTHRHD}\item [UART\_RX\_FLOW\_THRHD] 配置使用硬件流控时接收数据的最大值。 (R/W)
\label{fielddesc:UARTRXTOUTTHRHD}\item [UART\_RX\_TOUT\_THRHD] 配置接收器接收一个字节所需时间的阈值，单位是比特时间（即传输一个比特所需的时间）。
接收器接收一个字节所需时间超过阈值且 UART\_RX\_TOUT\_EN 置 1 时触发 UART\_RXFIFO\_TOUT\_INT 中断。 (R/W)
\label{fielddesc:UARTMEMFORCEPD}\item [UART\_MEM\_FORCE\_PD] 置位此位强制关闭 UART RAM。 (R/W)
\label{fielddesc:UARTMEMFORCEPU}\item [UART\_MEM\_FORCE\_PU] 置位此位强制开启 UART RAM。 (R/W)
\end{reglist}\end{regdesc}
\end{register}


\begin{register}{H}{UART\_INT\_RAW\_REG}{0x{}0004}\label{regdesc:UARTINTRAWREG}
\regfield{(reserved)}{12}{20}{000000000000}%
\regfield{UART\_WAKEUP\_INT\_RAW}{1}{19}{{0}}%
\regfield{UART\_AT\_CMD\_CHAR\_DET\_INT\_RAW}{1}{18}{{0}}%
\regfield{UART\_RS485\_CLASH\_INT\_RAW}{1}{17}{{0}}%
\regfield{UART\_RS485\_FRM\_ERR\_INT\_RAW}{1}{16}{{0}}%
\regfield{UART\_RS485\_PARITY\_ERR\_INT\_RAW}{1}{15}{{0}}%
\regfield{UART\_TX\_DONE\_INT\_RAW}{1}{14}{{0}}%
\regfield{UART\_TX\_BRK\_IDLE\_DONE\_INT\_RAW}{1}{13}{{0}}%
\regfield{UART\_TX\_BRK\_DONE\_INT\_RAW}{1}{12}{{0}}%
\regfield{UART\_GLITCH\_DET\_INT\_RAW}{1}{11}{{0}}%
\regfield{UART\_SW\_XOFF\_INT\_RAW}{1}{10}{{0}}%
\regfield{UART\_SW\_XON\_INT\_RAW}{1}{9}{{0}}%
\regfield{UART\_RXFIFO\_TOUT\_INT\_RAW}{1}{8}{{0}}%
\regfield{UART\_BRK\_DET\_INT\_RAW}{1}{7}{{0}}%
\regfield{UART\_CTS\_CHG\_INT\_RAW}{1}{6}{{0}}%
\regfield{UART\_DSR\_CHG\_INT\_RAW}{1}{5}{{0}}%
\regfield{UART\_RXFIFO\_OVF\_INT\_RAW}{1}{4}{{0}}%
\regfield{UART\_FRM\_ERR\_INT\_RAW}{1}{3}{{0}}%
\regfield{UART\_PARITY\_ERR\_INT\_RAW}{1}{2}{{0}}%
\regfield{UART\_TXFIFO\_EMPTY\_INT\_RAW}{1}{1}{{1}}%
\regfield{UART\_RXFIFO\_FULL\_INT\_RAW}{1}{0}{{0}}%
\reglabel{Reset}\regnewline%
\begin{regdesc}\begin{reglist}
\label{fielddesc:UARTRXFIFOFULLINTRAW}\item [UART\_RXFIFO\_FULL\_INT\_RAW] 接收器接收数据多于 UART\_RXFIFO\_FULL\_THRHD 的值时，该原始中断位翻转至高电平。 (R/WTC/SS)
\label{fielddesc:UARTTXFIFOEMPTYINTRAW}\item [UART\_TXFIFO\_EMPTY\_INT\_RAW] TX FIFO 中的数据少于 UART\_TXFIFO\_EMPTY\_THRHD 的值时，该原始中断位翻转至高电平。 (R/WTC/SS)
\label{fielddesc:UARTPARITYERRINTRAW}\item [UART\_PARITY\_ERR\_INT\_RAW] 接收器检测到数据奇偶检验位错误时，该原始中断位翻转至高电平。 (R/WTC/SS)
\label{fielddesc:UARTFRMERRINTRAW}\item [UART\_FRM\_ERR\_INT\_RAW] 接收器检测到数据帧错误时，该原始中断位翻转至高电平。 (R/WTC/SS)
\label{fielddesc:UARTRXFIFOOVFINTRAW}\item [UART\_RXFIFO\_OVF\_INT\_RAW] 接收器接收数据超过 RX FIFO 的存储容量时，该原始中断位翻转至高电平。 (R/WTC/SS)
\label{fielddesc:UARTDSRCHGINTRAW}\item [UART\_DSR\_CHG\_INT\_RAW] 接收器检测到 DSRn 信号的沿变化时，该原始中断位翻转至高电平。 (R/WTC/SS)
\label{fielddesc:UARTCTSCHGINTRAW}\item [UART\_CTS\_CHG\_INT\_RAW] 接收器检测到 CTSn 信号的沿变化时，该原始中断位翻转至高电平。 (R/WTC/SS)
\label{fielddesc:UARTBRKDETINTRAW}\item [UART\_BRK\_DET\_INT\_RAW] 接收器在停止位后检测到 0 时，该原始中断位翻转至高电平。 (R/WTC/SS)
\label{fielddesc:UARTRXFIFOTOUTINTRAW}\item [UART\_RXFIFO\_TOUT\_INT\_RAW] 接收器接收一个字节所需时间超过 UART\_RX\_TOUT\_THRHD 时，该原始中断位翻转至高电平。 (R/WTC/SS)
\label{fielddesc:UARTSWXONINTRAW}\item [UART\_SW\_XON\_INT\_RAW] 接收器接收到 XON 字符且 UART\_SW\_FLOW\_CON\_EN 置 1 时，该原始中断位翻转至高电平。 (R/WTC/SS)
\label{fielddesc:UARTSWXOFFINTRAW}\item [UART\_SW\_XOFF\_INT\_RAW] 接收器接收到 XOFF 字符且 UART\_SW\_FLOW\_CON\_EN 置 1 时，该原始中断位翻转至高电平。 (R/WTC/SS)
\label{fielddesc:UARTGLITCHDETINTRAW}\item [UART\_GLITCH\_DET\_INT\_RAW] 接收器在起始位的中点处检测到毛刺时，该原始中断位翻转至高电平。 (R/WTC/SS)

\item [见下页...]
\end{reglist}\end{regdesc}
\end{register}

\addtocounter{Regfloat}{-1}
\begin{register}{H}{UART\_INT\_RAW\_REG}{0x{}0004}
\begin{regdesc}\begin{reglist}
\item [接上页...]

\label{fielddesc:UARTTXBRKDONEINTRAW}\item [UART\_TX\_BRK\_DONE\_INT\_RAW] 发送器在发送完 TX FIFO 中所有数据后完成 NULL 字符的发送时，该原始中断位翻转至高电平。 (R/WTC/SS)
\label{fielddesc:UARTTXBRKIDLEDONEINTRAW}\item [UART\_TX\_BRK\_IDLE\_DONE\_INT\_RAW] 发送器发送完最后一个数据后的间隔时间达到阈值时，该原始中断位翻转至高电平。 (R/WTC/SS)
\label{fielddesc:UARTTXDONEINTRAW}\item [UART\_TX\_DONE\_INT\_RAW] 发送器发完 FIFO 中的所有数据后，该原始中断位翻转至高电平。 (R/WTC/SS)
\label{fielddesc:UARTRS485PARITYERRINTRAW}\item [UART\_RS485\_PARITY\_ERR\_INT\_RAW] RS485 模式下接收器检测到发送器回音的数据检验位错误时，该原始中断位翻转至高电平。 (R/WTC/SS)
\label{fielddesc:UARTRS485FRMERRINTRAW}\item [UART\_RS485\_FRM\_ERR\_INT\_RAW] RS485 模式下接收器检测到发送器回音的数据帧错误时，该原始中断位翻转至高电平。 (R/WTC/SS)
\label{fielddesc:UARTRS485CLASHINTRAW}\item [UART\_RS485\_CLASH\_INT\_RAW] RS485 模式下检测到发送器与接收器冲突时，该原始中断位翻转至高电平。 (R/WTC/SS)
\label{fielddesc:UARTATCMDCHARDETINTRAW}\item [UART\_AT\_CMD\_CHAR\_DET\_INT\_RAW] 接收器检测到配置的 UART\_AT\_CMD\_CHAR 时，该原始中断位翻转至高电平。 (R/WTC/SS)
\label{fielddesc:UARTWAKEUPINTRAW}\item [UART\_WAKEUP\_INT\_RAW] 输入 RXD 沿变化次数超过 Light-sleep 模式指定的 (UART\_ACTIVE\_THRESHOLD + 3) 值时，该原始中断位翻转至高电平。 (R/WTC/SS)
\end{reglist}\end{regdesc}
\end{register}


\begin{register}{H}{UART\_INT\_ST\_REG}{0x{}0008}\label{regdesc:UARTINTSTREG}
\regfield{(reserved)}{12}{20}{000000000000}%
\regfield{UART\_WAKEUP\_INT\_ST}{1}{19}{{0}}%
\regfield{UART\_AT\_CMD\_CHAR\_DET\_INT\_ST}{1}{18}{{0}}%
\regfield{UART\_RS485\_CLASH\_INT\_ST}{1}{17}{{0}}%
\regfield{UART\_RS485\_FRM\_ERR\_INT\_ST}{1}{16}{{0}}%
\regfield{UART\_RS485\_PARITY\_ERR\_INT\_ST}{1}{15}{{0}}%
\regfield{UART\_TX\_DONE\_INT\_ST}{1}{14}{{0}}%
\regfield{UART\_TX\_BRK\_IDLE\_DONE\_INT\_ST}{1}{13}{{0}}%
\regfield{UART\_TX\_BRK\_DONE\_INT\_ST}{1}{12}{{0}}%
\regfield{UART\_GLITCH\_DET\_INT\_ST}{1}{11}{{0}}%
\regfield{UART\_SW\_XOFF\_INT\_ST}{1}{10}{{0}}%
\regfield{UART\_SW\_XON\_INT\_ST}{1}{9}{{0}}%
\regfield{UART\_RXFIFO\_TOUT\_INT\_ST}{1}{8}{{0}}%
\regfield{UART\_BRK\_DET\_INT\_ST}{1}{7}{{0}}%
\regfield{UART\_CTS\_CHG\_INT\_ST}{1}{6}{{0}}%
\regfield{UART\_DSR\_CHG\_INT\_ST}{1}{5}{{0}}%
\regfield{UART\_RXFIFO\_OVF\_INT\_ST}{1}{4}{{0}}%
\regfield{UART\_FRM\_ERR\_INT\_ST}{1}{3}{{0}}%
\regfield{UART\_PARITY\_ERR\_INT\_ST}{1}{2}{{0}}%
\regfield{UART\_TXFIFO\_EMPTY\_INT\_ST}{1}{1}{{0}}%
\regfield{UART\_RXFIFO\_FULL\_INT\_ST}{1}{0}{{0}}%
\reglabel{Reset}\regnewline%
\begin{regdesc}\begin{reglist}
\label{fielddesc:UARTRXFIFOFULLINTST}\item [UART\_RXFIFO\_FULL\_INT\_ST] UART\_RXFIFO\_FULL\_INT\_ENA 置 1 时 UART\_RXFIFO\_FULL\_INT 的状态位。 (RO)
\label{fielddesc:UARTTXFIFOEMPTYINTST}\item [UART\_TXFIFO\_EMPTY\_INT\_ST] UART\_TXFIFO\_EMPTY\_INT\_ENA 置 1 时 UART\_TXFIFO\_EMPTY\_INT 的状态位。 (RO)
\label{fielddesc:UARTPARITYERRINTST}\item [UART\_PARITY\_ERR\_INT\_ST] UART\_PARITY\_ERR\_INT\_ENA 置 1 时 UART\_PARITY\_ERR\_INT 的状态位。 (RO)
\label{fielddesc:UARTFRMERRINTST}\item [UART\_FRM\_ERR\_INT\_ST] UART\_FRM\_ERR\_INT\_ENA 置 1 时 UART\_FRM\_ERR\_INT 的状态位。 (RO)
\label{fielddesc:UARTRXFIFOOVFINTST}\item [UART\_RXFIFO\_OVF\_INT\_ST] UART\_RXFIFO\_OVF\_INT\_ENA 置 1 时 UART\_RXFIFO\_OVF\_INT 的状态位。 (RO)
\label{fielddesc:UARTDSRCHGINTST}\item [UART\_DSR\_CHG\_INT\_ST] UART\_DSR\_CHG\_INT\_ENA 置 1 时 UART\_DSR\_CHG\_INT 的状态位。 (RO)
\label{fielddesc:UARTCTSCHGINTST}\item [UART\_CTS\_CHG\_INT\_ST] UART\_CTS\_CHG\_INT\_ENA 置 1 时 UART\_CTS\_CHG\_INT 的状态位。 (RO)
\label{fielddesc:UARTBRKDETINTST}\item [UART\_BRK\_DET\_INT\_ST] UART\_BRK\_DET\_INT\_ENA 置 1 时 UART\_BRK\_DET\_INT 的状态位。 (RO)
\label{fielddesc:UARTRXFIFOTOUTINTST}\item [UART\_RXFIFO\_TOUT\_INT\_ST] UART\_RXFIFO\_TOUT\_INT\_ENA 置 1 时 UART\_RXFIFO\_TOUT\_INT 的状态位。 (RO)
\label{fielddesc:UARTSWXONINTST}\item [UART\_SW\_XON\_INT\_ST] UART\_SW\_XON\_INT\_ENA 置 1 时 UART\_SW\_XON\_INT 的状态位。 (RO)
\label{fielddesc:UARTSWXOFFINTST}\item [UART\_SW\_XOFF\_INT\_ST] UART\_SW\_XOFF\_INT\_ENA 置 1 时 UART\_SW\_XOFF\_INT 的状态位。 (RO)
\label{fielddesc:UARTGLITCHDETINTST}\item [UART\_GLITCH\_DET\_INT\_ST] UART\_GLITCH\_DET\_INT\_ENA 置 1 时 UART\_GLITCH\_DET\_INT 的状态位。 (RO)

\item [见下页...]
\end{reglist}\end{regdesc}
\end{register}

\addtocounter{Regfloat}{-1}
\begin{register}{H}{UART\_INT\_ST\_REG}{0x{}0008}
\begin{regdesc}\begin{reglist}
\item [接上页...]

\label{fielddesc:UARTTXBRKDONEINTST}\item [UART\_TX\_BRK\_DONE\_INT\_ST] UART\_TX\_BRK\_DONE\_INT\_ENA 置 1 时 UART\_TX\_BRK\_DONE\_INT 的状态位。 (RO)
\label{fielddesc:UARTTXBRKIDLEDONEINTST}\item [UART\_TX\_BRK\_IDLE\_DONE\_INT\_ST] UART\_TX\_BRK\_IDLE\_DONE\_INT\_ENA 置 1 时 UART\_TX\_BRK\_IDLE\_DONE\_INT 的状态位。 (RO)
\label{fielddesc:UARTTXDONEINTST}\item [UART\_TX\_DONE\_INT\_ST] UART\_TX\_DONE\_INT\_ENA 置 1 时 UART\_TX\_DONE\_INT 的状态位。 (RO)
\label{fielddesc:UARTRS485PARITYERRINTST}\item [UART\_RS485\_PARITY\_ERR\_INT\_ST] UART\_RS485\_PARITY\_INT\_ENA 置 1 时 UART\_RS485\_PARITY\_ERR\_INT 的状态位。 (RO)
\label{fielddesc:UARTRS485FRMERRINTST}\item [UART\_RS485\_FRM\_ERR\_INT\_ST] UART\_RS485\_FRM\_ERR\_INT\_ENA 置 1 时 UART\_RS485\_FRM\_ERR\_INT 的状态位。 (RO)
\label{fielddesc:UARTRS485CLASHINTST}\item [UART\_RS485\_CLASH\_INT\_ST] UART\_RS485\_CLASH\_INT\_ENA 置 1 时 UART\_RS485\_CLASH\_INT 的状态位。 (RO)
\label{fielddesc:UARTATCMDCHARDETINTST}\item [UART\_AT\_CMD\_CHAR\_DET\_INT\_ST] UART\_AT\_CMD\_CHAR\_DET\_INT\_ENA 置 1 时 UART\_AT\_CMD\_CHAR\_DET\_INT 的状态位。 (RO)
\label{fielddesc:UARTWAKEUPINTST}\item [UART\_WAKEUP\_INT\_ST] UART\_WAKEUP\_INT\_ENA 置 1 时 UART\_WAKEUP\_INT 的状态位。 (RO)
\end{reglist}\end{regdesc}
\end{register}


\begin{register}{H}{UART\_INT\_ENA\_REG}{0x{}000C}\label{regdesc:UARTINTENAREG}
\regfield{(reserved)}{12}{20}{000000000000}%
\regfield{UART\_WAKEUP\_INT\_ENA}{1}{19}{{0}}%
\regfield{UART\_AT\_CMD\_CHAR\_DET\_INT\_ENA}{1}{18}{{0}}%
\regfield{UART\_RS485\_CLASH\_INT\_ENA}{1}{17}{{0}}%
\regfield{UART\_RS485\_FRM\_ERR\_INT\_ENA}{1}{16}{{0}}%
\regfield{UART\_RS485\_PARITY\_ERR\_INT\_ENA}{1}{15}{{0}}%
\regfield{UART\_TX\_DONE\_INT\_ENA}{1}{14}{{0}}%
\regfield{UART\_TX\_BRK\_IDLE\_DONE\_INT\_ENA}{1}{13}{{0}}%
\regfield{UART\_TX\_BRK\_DONE\_INT\_ENA}{1}{12}{{0}}%
\regfield{UART\_GLITCH\_DET\_INT\_ENA}{1}{11}{{0}}%
\regfield{UART\_SW\_XOFF\_INT\_ENA}{1}{10}{{0}}%
\regfield{UART\_SW\_XON\_INT\_ENA}{1}{9}{{0}}%
\regfield{UART\_RXFIFO\_TOUT\_INT\_ENA}{1}{8}{{0}}%
\regfield{UART\_BRK\_DET\_INT\_ENA}{1}{7}{{0}}%
\regfield{UART\_CTS\_CHG\_INT\_ENA}{1}{6}{{0}}%
\regfield{UART\_DSR\_CHG\_INT\_ENA}{1}{5}{{0}}%
\regfield{UART\_RXFIFO\_OVF\_INT\_ENA}{1}{4}{{0}}%
\regfield{UART\_FRM\_ERR\_INT\_ENA}{1}{3}{{0}}%
\regfield{UART\_PARITY\_ERR\_INT\_ENA}{1}{2}{{0}}%
\regfield{UART\_TXFIFO\_EMPTY\_INT\_ENA}{1}{1}{{0}}%
\regfield{UART\_RXFIFO\_FULL\_INT\_ENA}{1}{0}{{0}}%
\reglabel{Reset}\regnewline%
\begin{regdesc}\begin{reglist}
\label{fielddesc:UARTRXFIFOFULLINTENA}\item [UART\_RXFIFO\_FULL\_INT\_ENA] UART\_RXFIFO\_FULL\_INT 的使能位。 (R/W)
\label{fielddesc:UARTTXFIFOEMPTYINTENA}\item [UART\_TXFIFO\_EMPTY\_INT\_ENA] UART\_TXFIFO\_EMPTY\_INT 的使能位。 (R/W)
\label{fielddesc:UARTPARITYERRINTENA}\item [UART\_PARITY\_ERR\_INT\_ENA] UART\_PARITY\_ERR\_INT 的使能位。 (R/W)
\label{fielddesc:UARTFRMERRINTENA}\item [UART\_FRM\_ERR\_INT\_ENA] UART\_FRM\_ERR\_INT 的使能位。 (R/W)
\label{fielddesc:UARTRXFIFOOVFINTENA}\item [UART\_RXFIFO\_OVF\_INT\_ENA] UART\_RXFIFO\_OVF\_INT 的使能位。 (R/W)
\label{fielddesc:UARTDSRCHGINTENA}\item [UART\_DSR\_CHG\_INT\_ENA] UART\_DSR\_CHG\_INT 的使能位。 (R/W)
\label{fielddesc:UARTCTSCHGINTENA}\item [UART\_CTS\_CHG\_INT\_ENA] UART\_CTS\_CHG\_INT 的使能位。 (R/W)
\label{fielddesc:UARTBRKDETINTENA}\item [UART\_BRK\_DET\_INT\_ENA] UART\_BRK\_DET\_INT 的使能位。 (R/W)
\label{fielddesc:UARTRXFIFOTOUTINTENA}\item [UART\_RXFIFO\_TOUT\_INT\_ENA] UART\_RXFIFO\_TOUT\_INT 的使能位。 (R/W)
\label{fielddesc:UARTSWXONINTENA}\item [UART\_SW\_XON\_INT\_ENA] UART\_SW\_XON\_INT 的使能位。 (R/W)
\label{fielddesc:UARTSWXOFFINTENA}\item [UART\_SW\_XOFF\_INT\_ENA] UART\_SW\_XOFF\_INT 的使能位。 (R/W)
\label{fielddesc:UARTGLITCHDETINTENA}\item [UART\_GLITCH\_DET\_INT\_ENA] UART\_GLITCH\_DET\_INT 的使能位。 (R/W)
\label{fielddesc:UARTTXBRKDONEINTENA}\item [UART\_TX\_BRK\_DONE\_INT\_ENA] UART\_TX\_BRK\_DONE\_INT 的使能位。 (R/W)
\label{fielddesc:UARTTXBRKIDLEDONEINTENA}\item [UART\_TX\_BRK\_IDLE\_DONE\_INT\_ENA] UART\_TX\_BRK\_IDLE\_DONE\_INT 的使能位。 (R/W)
\label{fielddesc:UARTTXDONEINTENA}\item [UART\_TX\_DONE\_INT\_ENA] UART\_TX\_DONE\_INT 的使能位。 (R/W)
\label{fielddesc:UARTRS485PARITYERRINTENA}\item [UART\_RS485\_PARITY\_ERR\_INT\_ENA] UART\_RS485\_PARITY\_ERR\_INT 的使能位。 (R/W)
\label{fielddesc:UARTRS485FRMERRINTENA}\item [UART\_RS485\_FRM\_ERR\_INT\_ENA] UART\_RS485\_PARITY\_ERR\_INT 的使能位。 (R/W)
\label{fielddesc:UARTRS485CLASHINTENA}\item [UART\_RS485\_CLASH\_INT\_ENA] UART\_RS485\_CLASH\_INT 的使能位。 (R/W)
\label{fielddesc:UARTATCMDCHARDETINTENA}\item [UART\_AT\_CMD\_CHAR\_DET\_INT\_ENA] UART\_AT\_CMD\_CHAR\_DET\_INT 的使能位。 (R/W)
\label{fielddesc:UARTWAKEUPINTENA}\item [UART\_WAKEUP\_INT\_ENA] UART\_WAKEUP\_INT 的使能位。 (R/W)
\end{reglist}\end{regdesc}
\end{register}


\begin{register}{H}{UART\_INT\_CLR\_REG}{0x{}0010}\label{regdesc:UARTINTCLRREG}
\regfield{(reserved)}{12}{20}{000000000000}%
\regfield{UART\_WAKEUP\_INT\_CLR}{1}{19}{{0}}%
\regfield{UART\_AT\_CMD\_CHAR\_DET\_INT\_CLR}{1}{18}{{0}}%
\regfield{UART\_RS485\_CLASH\_INT\_CLR}{1}{17}{{0}}%
\regfield{UART\_RS485\_FRM\_ERR\_INT\_CLR}{1}{16}{{0}}%
\regfield{UART\_RS485\_PARITY\_ERR\_INT\_CLR}{1}{15}{{0}}%
\regfield{UART\_TX\_DONE\_INT\_CLR}{1}{14}{{0}}%
\regfield{UART\_TX\_BRK\_IDLE\_DONE\_INT\_CLR}{1}{13}{{0}}%
\regfield{UART\_TX\_BRK\_DONE\_INT\_CLR}{1}{12}{{0}}%
\regfield{UART\_GLITCH\_DET\_INT\_CLR}{1}{11}{{0}}%
\regfield{UART\_SW\_XOFF\_INT\_CLR}{1}{10}{{0}}%
\regfield{UART\_SW\_XON\_INT\_CLR}{1}{9}{{0}}%
\regfield{UART\_RXFIFO\_TOUT\_INT\_CLR}{1}{8}{{0}}%
\regfield{UART\_BRK\_DET\_INT\_CLR}{1}{7}{{0}}%
\regfield{UART\_CTS\_CHG\_INT\_CLR}{1}{6}{{0}}%
\regfield{UART\_DSR\_CHG\_INT\_CLR}{1}{5}{{0}}%
\regfield{UART\_RXFIFO\_OVF\_INT\_CLR}{1}{4}{{0}}%
\regfield{UART\_FRM\_ERR\_INT\_CLR}{1}{3}{{0}}%
\regfield{UART\_PARITY\_ERR\_INT\_CLR}{1}{2}{{0}}%
\regfield{UART\_TXFIFO\_EMPTY\_INT\_CLR}{1}{1}{{0}}%
\regfield{UART\_RXFIFO\_FULL\_INT\_CLR}{1}{0}{{0}}%
\reglabel{Reset}\regnewline%
\begin{regdesc}\begin{reglist}
\label{fielddesc:UARTRXFIFOFULLINTCLR}\item [UART\_RXFIFO\_FULL\_INT\_CLR] 置位此位清除 UART\_RXFIFO\_FULL\_INT 中断。 (WT)
\label{fielddesc:UARTTXFIFOEMPTYINTCLR}\item [UART\_TXFIFO\_EMPTY\_INT\_CLR] 置位此位清除 UART\_TXFIFO\_EMPTY\_INT 中断。 (WT)
\label{fielddesc:UARTPARITYERRINTCLR}\item [UART\_PARITY\_ERR\_INT\_CLR] 置位此位清除 UART\_PARITY\_ERR\_INT 中断。 (WT)
\label{fielddesc:UARTFRMERRINTCLR}\item [UART\_FRM\_ERR\_INT\_CLR] 置位此位清除 UART\_FRM\_ERR\_INT 中断。 (WT)
\label{fielddesc:UARTRXFIFOOVFINTCLR}\item [UART\_RXFIFO\_OVF\_INT\_CLR] 置位此位清除 UART\_RXFIFO\_OVF\_INT 中断。 (WT)
\label{fielddesc:UARTDSRCHGINTCLR}\item [UART\_DSR\_CHG\_INT\_CLR] 置位此位清除 UART\_DSR\_CHG\_INT 中断。 (WT)
\label{fielddesc:UARTCTSCHGINTCLR}\item [UART\_CTS\_CHG\_INT\_CLR] 置位此位清除 UART\_CTS\_CHG\_INT 中断。 (WT)
\label{fielddesc:UARTBRKDETINTCLR}\item [UART\_BRK\_DET\_INT\_CLR] 置位此位清除 UART\_BRK\_DET\_INT 中断。 (WT)
\label{fielddesc:UARTRXFIFOTOUTINTCLR}\item [UART\_RXFIFO\_TOUT\_INT\_CLR] 置位此位清除 UART\_RXFIFO\_TOUT\_INT 中断。 (WT)
\label{fielddesc:UARTSWXONINTCLR}\item [UART\_SW\_XON\_INT\_CLR] 置位此位清除 UART\_SW\_XON\_INT 中断。 (WT)
\label{fielddesc:UARTSWXOFFINTCLR}\item [UART\_SW\_XOFF\_INT\_CLR] 置位此位清除 UART\_SW\_XOFF\_INT 中断。 (WT)
\label{fielddesc:UARTGLITCHDETINTCLR}\item [UART\_GLITCH\_DET\_INT\_CLR] 置位此位清除 UART\_GLITCH\_DET\_INT 中断。 (WT)
\label{fielddesc:UARTTXBRKDONEINTCLR}\item [UART\_TX\_BRK\_DONE\_INT\_CLR] 置位此位清除 UART\_TX\_BRK\_DONE\_INT 中断。 (WT)
\label{fielddesc:UARTTXBRKIDLEDONEINTCLR}\item [UART\_TX\_BRK\_IDLE\_DONE\_INT\_CLR] 置位此位清除 UART\_TX\_BRK\_IDLE\_DONE\_INT 中断。 (WT)
\label{fielddesc:UARTTXDONEINTCLR}\item [UART\_TX\_DONE\_INT\_CLR] 置位此位清除 UART\_TX\_DONE\_INT 中断。 (WT)
\label{fielddesc:UARTRS485PARITYERRINTCLR}\item [UART\_RS485\_PARITY\_ERR\_INT\_CLR] 置位此位清除 UART\_RS485\_PARITY\_ERR\_INT 中断。 (WT)
\label{fielddesc:UARTRS485FRMERRINTCLR}\item [UART\_RS485\_FRM\_ERR\_INT\_CLR] 置位此位清除 UART\_RS485\_FRM\_ERR\_INT 中断。 (WT)
\label{fielddesc:UARTRS485CLASHINTCLR}\item [UART\_RS485\_CLASH\_INT\_CLR] 置位此位清除 UART\_RS485\_CLASH\_INT 中断。 (WT)
\label{fielddesc:UARTATCMDCHARDETINTCLR}\item [UART\_AT\_CMD\_CHAR\_DET\_INT\_CLR] 置位此位清除 UART\_AT\_CMD\_CHAR\_DET\_INT 中断。 (WT)
\label{fielddesc:UARTWAKEUPINTCLR}\item [UART\_WAKEUP\_INT\_CLR] 置位此位清除 UART\_WAKEUP\_INT 中断。 (WT)
\end{reglist}\end{regdesc}
\end{register}


\begin{register}{H}{UART\_CLKDIV\_REG}{0x{}0014}\label{regdesc:UARTCLKDIVREG}
\regfield{(reserved)}{8}{24}{00000000}%
\regfield{UART\_CLKDIV\_FRAG}{4}{20}{{0x{}0}}%
\regfield{(reserved)}{8}{12}{00000000}%
\regfield{UART\_CLKDIV}{12}{0}{{0x{}2b6}}%
\reglabel{Reset}\regnewline%
\begin{regdesc}\begin{reglist}
\label{fielddesc:UARTCLKDIV}\item [UART\_CLKDIV] 分频系数的整数部分。 (R/W)
\label{fielddesc:UARTCLKDIVFRAG}\item [UART\_CLKDIV\_FRAG] 分频系数的小数部分。 (R/W)
\end{reglist}\end{regdesc}
\end{register}


\begin{register}{H}{UART\_RX\_FILT\_REG}{0x{}0018}\label{regdesc:UARTRXFILTREG}
\regfield{(reserved)}{23}{9}{00000000000000000000000}%
\regfield{UART\_GLITCH\_FILT\_EN}{1}{8}{{0}}%
\regfield{UART\_GLITCH\_FILT}{8}{0}{{0x{}8}}%
\reglabel{Reset}\regnewline%
\begin{regdesc}\begin{reglist}
\label{fielddesc:UARTGLITCHFILT}\item [UART\_GLITCH\_FILT] 宽度小于该字段值的输入脉冲会被忽略。 (R/W)
\label{fielddesc:UARTGLITCHFILTEN}\item [UART\_GLITCH\_FILT\_EN] 置位此位，使能 RX 信号滤波器。 (R/W)
\end{reglist}\end{regdesc}
\end{register}


\begin{register}{H}{UART\_CONF0\_REG}{0x{}0020}\label{regdesc:UARTCONF0REG}
\regfield{(reserved)}{3}{29}{000}%
\regfield{UART\_MEM\_CLK\_EN}{1}{28}{{1}}%
\regfield{UART\_AUTOBAUD\_EN}{1}{27}{{0}}%
\regfield{UART\_ERR\_WR\_MASK}{1}{26}{{0}}%
\regfield{UART\_CLK\_EN}{1}{25}{{0}}%
\regfield{UART\_DTR\_INV}{1}{24}{{0}}%
\regfield{UART\_RTS\_INV}{1}{23}{{0}}%
\regfield{UART\_TXD\_INV}{1}{22}{{0}}%
\regfield{UART\_DSR\_INV}{1}{21}{{0}}%
\regfield{UART\_CTS\_INV}{1}{20}{{0}}%
\regfield{UART\_RXD\_INV}{1}{19}{{0}}%
\regfield{UART\_TXFIFO\_RST}{1}{18}{{0}}%
\regfield{UART\_RXFIFO\_RST}{1}{17}{{0}}%
\regfield{UART\_IRDA\_EN}{1}{16}{{0}}%
\regfield{UART\_TX\_FLOW\_EN}{1}{15}{{0}}%
\regfield{UART\_LOOPBACK}{1}{14}{{0}}%
\regfield{UART\_IRDA\_RX\_INV}{1}{13}{{0}}%
\regfield{UART\_IRDA\_TX\_INV}{1}{12}{{0}}%
\regfield{UART\_IRDA\_WCTL}{1}{11}{{0}}%
\regfield{UART\_IRDA\_TX\_EN}{1}{10}{{0}}%
\regfield{UART\_IRDA\_DPLX}{1}{9}{{0}}%
\regfield{UART\_TXD\_BRK}{1}{8}{{0}}%
\regfield{UART\_SW\_DTR}{1}{7}{{0}}%
\regfield{UART\_SW\_RTS}{1}{6}{{0}}%
\regfield{UART\_STOP\_BIT\_NUM}{2}{4}{{1}}%
\regfield{UART\_BIT\_NUM}{2}{2}{{3}}%
\regfield{UART\_PARITY\_EN}{1}{1}{{0}}%
\regfield{UART\_PARITY}{1}{0}{{0}}%
\reglabel{Reset}\regnewline%
\begin{regdesc}\begin{reglist}
\label{fielddesc:UARTPARITY}\item [UART\_PARITY] 配置奇偶检验方式。 (R/W)
\label{fielddesc:UARTPARITYEN}\item [UART\_PARITY\_EN] 置位此位使能 UART 奇偶检验。 (R/W)
\label{fielddesc:UARTBITNUM}\item [UART\_BIT\_NUM] 设置数据长度。 (R/W)
\label{fielddesc:UARTSTOPBITNUM}\item [UART\_STOP\_BIT\_NUM] 设置停止位的长度。 (R/W)
\label{fielddesc:UARTSWRTS}\item [UART\_SW\_RTS] 该位用于配置软件流控使用的软件 RTS 信号。 (R/W)
\label{fielddesc:UARTSWDTR}\item [UART\_SW\_DTR] 该位用于配置软件流控使用的软件 DTR 信号。 (R/W)
\label{fielddesc:UARTTXDBRK}\item [UART\_TXD\_BRK] 置位此位，使能发送器在发完数据后发送 NULL。 (R/W)
\label{fielddesc:UARTIRDADPLX}\item [UART\_IRDA\_DPLX] 置位此位开启 IrDA 回环测试模式。 (R/W)
\label{fielddesc:UARTIRDATXEN}\item [UART\_IRDA\_TX\_EN] IrDA 发送器的启动使能位。 (R/W)
\label{fielddesc:UARTIRDAWCTL}\item [UART\_IRDA\_WCTL] 0：将 IrDA 发送器的第 11 位置 0；1：IrDA 发送器的第 11 位与第 10 位相同。 (R/W)
\label{fielddesc:UARTIRDATXINV}\item [UART\_IRDA\_TX\_INV] 置位此位翻转 IrDA 发送器的电平。 (R/W)
\label{fielddesc:UARTIRDARXINV}\item [UART\_IRDA\_RX\_INV] 置位此位翻转 IrDA 接收器的电平。 (R/W)
\label{fielddesc:UARTLOOPBACK}\item [UART\_LOOPBACK] 置位此位开启 UART 回环测试模式。 (R/W)
\label{fielddesc:UARTTXFLOWEN}\item [UART\_TX\_FLOW\_EN] 置位此位使能发送器的流控功能。 (R/W)
\label{fielddesc:UARTIRDAEN}\item [UART\_IRDA\_EN] 置位此位使能 IrDA 协议。 (R/W)
\label{fielddesc:UARTRXFIFORST}\item [UART\_RXFIFO\_RST] 置位此位复位 UART RX FIFO。 (R/W)
\label{fielddesc:UARTTXFIFORST}\item [UART\_TXFIFO\_RST] 置位此位复位 UART TX FIFO。 (R/W)
\label{fielddesc:UARTRXDINV}\item [UART\_RXD\_INV] 置位此位翻转 UART RXD 信号电平。 (R/W)
\label{fielddesc:UARTCTSINV}\item [UART\_CTS\_INV] 置位此位翻转 UART CTS 信号电平。 (R/W)
\label{fielddesc:UARTDSRINV}\item [UART\_DSR\_INV] 置位此位翻转 UART DSR 信号电平。 (R/W)
\label{fielddesc:UARTTXDINV}\item [UART\_TXD\_INV] 置位此位翻转 UART TXD 信号电平。 (R/W)
\label{fielddesc:UARTRTSINV}\item [UART\_RTS\_INV] 置位此位翻转 UART RTS 信号电平。 (R/W)
\label{fielddesc:UARTDTRINV}\item [UART\_DTR\_INV] 置位此位翻转 UART DTR 信号电平。 (R/W)

\item [见下页...]
\end{reglist}\end{regdesc}
\end{register}

\addtocounter{Regfloat}{-1}
\begin{register}{H}{UART\_CONF0\_REG}{0x{}0020}
\begin{regdesc}\begin{reglist}
\item [接上页...]

\label{fielddesc:UARTCLKEN}\item [UART\_CLK\_EN] 0：仅在应用写寄存器时支持时钟；1：强制为寄存器开启时钟。 (R/W)
\label{fielddesc:UARTERRWRMASK}\item [UART\_ERR\_WR\_MASK] 0：若数据错误，接收器仍存储；1：若数据错误，接收器不再将数据存入 FIFO。 (R/W)
\label{fielddesc:UARTAUTOBAUDEN}\item [UART\_AUTOBAUD\_EN] 波特率检测的使能信号。 (R/W)
\label{fielddesc:UARTMEMCLKEN}\item [UART\_MEM\_CLK\_EN] UART RAM 门控使能信号。 (R/W)
\end{reglist}\end{regdesc}
\end{register}


\begin{register}{H}{UART\_CONF1\_REG}{0x{}0024}\label{regdesc:UARTCONF1REG}
\regfield{(reserved)}{8}{24}{00000000}%
\regfield{UART\_RX\_TOUT\_EN}{1}{23}{{0}}%
\regfield{UART\_RX\_FLOW\_EN}{1}{22}{{0}}%
\regfield{UART\_RX\_TOUT\_FLOW\_DIS}{1}{21}{{0}}%
\regfield{UART\_DIS\_RX\_DAT\_OVF}{1}{20}{{0}}%
\regfield{UART\_TXFIFO\_EMPTY\_THRHD}{10}{10}{{0x{}60}}%
\regfield{UART\_RXFIFO\_FULL\_THRHD}{10}{0}{{0x{}60}}%
\reglabel{Reset}\regnewline%
\begin{regdesc}\begin{reglist}
\label{fielddesc:UARTRXFIFOFULLTHRHD}\item [UART\_RXFIFO\_FULL\_THRHD] 接收器接收数据多于该字段的值时产生 UART\_RXFIFO\_FULL\_INT 中断。 (R/W)
\label{fielddesc:UARTTXFIFOEMPTYTHRHD}\item [UART\_TXFIFO\_EMPTY\_THRHD] TX FIFO 中的数据少于该字段的值时产生 UART\_TXFIFO\_EMPTY\_INT 中断。 (R/W)
\label{fielddesc:UARTDISRXDATOVF}\item [UART\_DIS\_RX\_DAT\_OVF] 关闭 UART RX 数据溢出检测。 (R/W)
\label{fielddesc:UARTRXTOUTFLOWDIS}\item [UART\_RX\_TOUT\_FLOW\_DIS] 使用硬件流控时置位此位停止堆积 idle\_cnt。 (R/W)
\label{fielddesc:UARTRXFLOWEN}\item [UART\_RX\_FLOW\_EN] UART 接收器流控功能的使能位。 (R/W)
\label{fielddesc:UARTRXTOUTEN}\item [UART\_RX\_TOUT\_EN] UART 接收器超时功能的使能位。 (R/W)
\end{reglist}\end{regdesc}
\end{register}


\begin{register}{H}{UART\_FLOW\_CONF\_REG}{0x{}0034}\label{regdesc:UARTFLOWCONFREG}
\regfield{(reserved)}{26}{6}{00000000000000000000000000}%
\regfield{UART\_SEND\_XOFF}{1}{5}{{0}}%
\regfield{UART\_SEND\_XON}{1}{4}{{0}}%
\regfield{UART\_FORCE\_XOFF}{1}{3}{{0}}%
\regfield{UART\_FORCE\_XON}{1}{2}{{0}}%
\regfield{UART\_XONOFF\_DEL}{1}{1}{{0}}%
\regfield{UART\_SW\_FLOW\_CON\_EN}{1}{0}{{0}}%
\reglabel{Reset}\regnewline%
\begin{regdesc}\begin{reglist}
\label{fielddesc:UARTSWFLOWCONEN}\item [UART\_SW\_FLOW\_CON\_EN] 置位此位使能软件流控。UART 接收到 UART\_XON\_CHAR 或 UART\_XOFF\_CHAR 配置的流控字符 XON 或 XOFF 时，UART\_SW\_XON\_INT 或 UART\_SW\_XOFF\_INT 中断可在使能时触发。 (R/W)
\label{fielddesc:UARTXONOFFDEL}\item [UART\_XONOFF\_DEL] 置位此位移除接收数据中的流控字符。 (R/W)
\label{fielddesc:UARTFORCEXON}\item [UART\_FORCE\_XON] 置位此位让发送器继续发送数据。 (R/W)
\label{fielddesc:UARTFORCEXOFF}\item [UART\_FORCE\_XOFF] 置位此位让发送器停止发送数据。 (R/W)
\label{fielddesc:UARTSENDXON}\item [UART\_SEND\_XON] 置位此位发送 XON 字符。 此位由硬件自动清除。 (R/W/SS/SC)
\label{fielddesc:UARTSENDXOFF}\item [UART\_SEND\_XOFF] 置位此位发送 XOFF 字符。 此位由硬件自动清除。 (R/W/SS/SC)
\end{reglist}\end{regdesc}
\end{register}


\begin{register}{H}{UART\_SLEEP\_CONF\_REG}{0x{}0038}\label{regdesc:UARTSLEEPCONFREG}
\regfield{(reserved)}{22}{10}{0000000000000000000000}%
\regfield{UART\_ACTIVE\_THRESHOLD}{10}{0}{{0x{}f0}}%
\reglabel{Reset}\regnewline%
\begin{regdesc}\begin{reglist}
\label{fielddesc:UARTACTIVETHRESHOLD}\item [UART\_ACTIVE\_THRESHOLD] 输入 RXD 沿变化次数超过该字段的值加 3 时，UART 从 Light-sleep 模式唤醒。 (R/W)
\end{reglist}\end{regdesc}
\end{register}


\begin{register}{H}{UART\_SWFC\_CONF0\_REG}{0x{}003C}\label{regdesc:UARTSWFCCONF0REG}
\regfield{(reserved)}{14}{18}{00000000000000}%
\regfield{UART\_XOFF\_CHAR}{8}{10}{{0x{}13}}%
\regfield{UART\_XOFF\_THRESHOLD}{10}{0}{{0x{}e0}}%
\reglabel{Reset}\regnewline%
\begin{regdesc}\begin{reglist}
\label{fielddesc:UARTXOFFTHRESHOLD}\item [UART\_XOFF\_THRESHOLD] RX FIFO 中的数据超过该字段的值且 UART\_SW\_FLOW\_CON\_EN 置 1 时，发送 XOFF 字符。 (R/W)
\label{fielddesc:UARTXOFFCHAR}\item [UART\_XOFF\_CHAR] 存储 XOFF 流控字符。 (R/W)
\end{reglist}\end{regdesc}
\end{register}


\begin{register}{H}{UART\_SWFC\_CONF1\_REG}{0x{}0040}\label{regdesc:UARTSWFCCONF1REG}
\regfield{(reserved)}{14}{18}{00000000000000}%
\regfield{UART\_XON\_CHAR}{8}{10}{{0x{}11}}%
\regfield{UART\_XON\_THRESHOLD}{10}{0}{{0x{}0}}%
\reglabel{Reset}\regnewline%
\begin{regdesc}\begin{reglist}
\label{fielddesc:UARTXONTHRESHOLD}\item [UART\_XON\_THRESHOLD] RX FIFO 中的数据小于该字段的值且 UART\_SW\_FLOW\_CON\_EN 置 1 时，发送 XON 字符。 (R/W)
\label{fielddesc:UARTXONCHAR}\item [UART\_XON\_CHAR] 存储 XON 流控字符。 (R/W)
\end{reglist}\end{regdesc}
\end{register}


\begin{register}{H}{UART\_TXBRK\_CONF\_REG}{0x{}0044}\label{regdesc:UARTTXBRKCONFREG}
\regfield{(reserved)}{24}{8}{000000000000000000000000}%
\regfield{UART\_TX\_BRK\_NUM}{8}{0}{{0x{}a}}%
\reglabel{Reset}\regnewline%
\begin{regdesc}\begin{reglist}
\label{fielddesc:UARTTXBRKNUM}\item [UART\_TX\_BRK\_NUM] 配置数据发完后待发 NULL 字符的数量。UART\_TXD\_BRK 置 1 时有意义。 (R/W)
\end{reglist}\end{regdesc}
\end{register}


\begin{register}{H}{UART\_IDLE\_CONF\_REG}{0x{}0048}\label{regdesc:UARTIDLECONFREG}
\regfield{(reserved)}{12}{20}{000000000000}%
\regfield{UART\_TX\_IDLE\_NUM}{10}{10}{{0x{}100}}%
\regfield{UART\_RX\_IDLE\_THRHD}{10}{0}{{0x{}100}}%
\reglabel{Reset}\regnewline%
\begin{regdesc}\begin{reglist}
\label{fielddesc:UARTRXIDLETHRHD}\item [UART\_RX\_IDLE\_THRHD] 接收器接收一字节数据所需时间超过该字段的值时产生帧结束信号，单位是比特时间（即传输一个比特所需的时间）。 (R/W)
\label{fielddesc:UARTTXIDLENUM}\item [UART\_TX\_IDLE\_NUM] 配置两次数据传输的间隔时间，单位是比特时间（即传输一个比特所需的时间）。 (R/W)
\end{reglist}\end{regdesc}
\end{register}


\begin{register}{H}{UART\_RS485\_CONF\_REG}{0x{}004C}\label{regdesc:UARTRS485CONFREG}
\regfield{(reserved)}{22}{10}{0000000000000000000000}%
\regfield{UART\_RS485\_TX\_DLY\_NUM}{4}{6}{{0}}%
\regfield{UART\_RS485\_RX\_DLY\_NUM}{1}{5}{{0}}%
\regfield{UART\_RS485RXBY\_TX\_EN}{1}{4}{{0}}%
\regfield{UART\_RS485TX\_RX\_EN}{1}{3}{{0}}%
\regfield{UART\_DL1\_EN}{1}{2}{{0}}%
\regfield{UART\_DL0\_EN}{1}{1}{{0}}%
\regfield{UART\_RS485\_EN}{1}{0}{{0}}%
\reglabel{Reset}\regnewline%
\begin{regdesc}\begin{reglist}
\label{fielddesc:UARTRS485EN}\item [UART\_RS485\_EN] 置位此位选择 RS485 模式。 (R/W)
\label{fielddesc:UARTDL0EN}\item [UART\_DL0\_EN] 配置是否在起始位之前增加一位延时。\\
0：不增加\
1：增加\\ (R/W)
\label{fielddesc:UARTDL1EN}\item [UART\_DL1\_EN] 配置是否在停止位之后增加一位延时。\\
0：不增加\\
1：增加\\ (R/W)
\label{fielddesc:UARTRS485TXRXEN}\item [UART\_RS485TX\_RX\_EN] 发送器在 RS485 模式下发送数据时，置位此位使能接收器接收数据。 (R/W)
\label{fielddesc:UARTRS485RXBYTXEN}\item [UART\_RS485RXBY\_TX\_EN] 1'h1： RS485 接收器线路繁忙时使能 RS485 发送器发送数据。 (R/W)
\label{fielddesc:UARTRS485RXDLYNUM}\item [UART\_RS485\_RX\_DLY\_NUM] 延迟接收器的内部数据信号。 (R/W)
\label{fielddesc:UARTRS485TXDLYNUM}\item [UART\_RS485\_TX\_DLY\_NUM] 延迟发送器的内部数据信号。 (R/W)
\end{reglist}\end{regdesc}
\end{register}


\begin{register}{H}{UART\_CLK\_CONF\_REG}{0x{}0078}\label{regdesc:UARTCLKCONFREG}
\regfield{(reserved)}{4}{28}{0000}%
\regfield{UART\_RX\_RST\_CORE}{1}{27}{{0}}%
\regfield{UART\_TX\_RST\_CORE}{1}{26}{{0}}%
\regfield{UART\_RX\_SCLK\_EN}{1}{25}{{1}}%
\regfield{UART\_TX\_SCLK\_EN}{1}{24}{{1}}%
\regfield{UART\_RST\_CORE}{1}{23}{{0}}%
\regfield{UART\_SCLK\_EN}{1}{22}{{1}}%
\regfield{UART\_SCLK\_SEL}{2}{20}{{3}}%
\regfield{UART\_SCLK\_DIV\_NUM}{8}{12}{{0x{}1}}%
\regfield{UART\_SCLK\_DIV\_A}{6}{6}{{0x{}0}}%
\regfield{UART\_SCLK\_DIV\_B}{6}{0}{{0x{}0}}%
\reglabel{Reset}\regnewline%
\begin{regdesc}\begin{reglist}
\label{fielddesc:UARTSCLKDIVB}\item [UART\_SCLK\_DIV\_B] 分频系数的分母。 (R/W)
\label{fielddesc:UARTSCLKDIVA}\item [UART\_SCLK\_DIV\_A] 分频系数的分子。 (R/W)
\label{fielddesc:UARTSCLKDIVNUM}\item [UART\_SCLK\_DIV\_NUM] 分频系数的整数部分。 (R/W)
\label{fielddesc:UARTSCLKSEL}\item [UART\_SCLK\_SEL] 选择 UART 时钟源。1：APB\_CLK；2：RC\_FAST\_CLK；3： XTAL\_CLK。 (R/W)
\label{fielddesc:UARTSCLKEN}\item [UART\_SCLK\_EN] 置位此位，使能 UART TX/RX 使能。 (R/W)
\label{fielddesc:UARTRSTCORE}\item [UART\_RST\_CORE] 向此位先写 1 后写 0，复位 UART TX/RX。 (R/W)
\label{fielddesc:UARTTXSCLKEN}\item [UART\_TX\_SCLK\_EN] 置位此位，使能 UART TX 时钟。 (R/W)
\label{fielddesc:UARTRXSCLKEN}\item [UART\_RX\_SCLK\_EN] 置位此位，使能 UART RX 时钟。 (R/W)
\label{fielddesc:UARTTXRSTCORE}\item [UART\_TX\_RST\_CORE] 向此位先写 1 后写 0，复位 UART TX。 (R/W)
\label{fielddesc:UARTRXRSTCORE}\item [UART\_RX\_RST\_CORE] 向此位先写 1 后写 0，复位 UART RX。 (R/W)
\end{reglist}\end{regdesc}
\end{register}


\begin{register}{H}{UART\_STATUS\_REG}{0x{}001C}\label{regdesc:UARTSTATUSREG}
\regfield{UART\_TXD}{1}{31}{{1}}%
\regfield{UART\_RTSN}{1}{30}{{1}}%
\regfield{UART\_DTRN}{1}{29}{{1}}%
\regfield{(reserved)}{3}{26}{000}%
\regfield{UART\_TXFIFO\_CNT}{10}{16}{{0}}%
\regfield{UART\_RXD}{1}{15}{{1}}%
\regfield{UART\_CTSN}{1}{14}{{1}}%
\regfield{UART\_DSRN}{1}{13}{{0}}%
\regfield{(reserved)}{3}{10}{000}%
\regfield{UART\_RXFIFO\_CNT}{10}{0}{{0}}%
\reglabel{Reset}\regnewline%
\begin{regdesc}\begin{reglist}
\label{fielddesc:UARTRXFIFOCNT}\item [UART\_RXFIFO\_CNT] 存储 RX FIFO 中有效数据的字节数。 (RO)
\label{fielddesc:UARTDSRN}\item [UART\_DSRN] 该位表示内部 UART DSR 信号的电平值。 (RO)
\label{fielddesc:UARTCTSN}\item [UART\_CTSN] 该位表示内部 UART CTS 信号的电平值。 (RO)
\label{fielddesc:UARTRXD}\item [UART\_RXD] 该位表示内部 UART RXD 信号的电平值。 (RO)
\label{fielddesc:UARTTXFIFOCNT}\item [UART\_TXFIFO\_CNT] 存储 TX FIFO 中数据的字节数。 (RO)
\label{fielddesc:UARTDTRN}\item [UART\_DTRN] 此位表示内部 UART DTR 信号的电平。 (RO)
\label{fielddesc:UARTRTSN}\item [UART\_RTSN] 此位表示内部 UART RTS 信号的电平。 (RO)
\label{fielddesc:UARTTXD}\item [UART\_TXD] 此位表示内部 UART TXD 信号的电平。 (RO)
\end{reglist}\end{regdesc}
\end{register}


\begin{register}{H}{UART\_MEM\_TX\_STATUS\_REG}{0x{}0064}\label{regdesc:UARTMEMTXSTATUSREG}
\regfield{(reserved)}{11}{21}{00000000000}%
\regfield{UART\_TX\_RADDR}{10}{11}{{0x{}0}}%
\regfield{(reserved)}{1}{10}{0}%
\regfield{UART\_APB\_TX\_WADDR}{10}{0}{{0x{}0}}%
\reglabel{Reset}\regnewline%
\begin{regdesc}\begin{reglist}
\label{fielddesc:UARTAPBTXWADDR}\item [UART\_APB\_TX\_WADDR] 在软件通过 APB 总线写 TX FIFO 时存储 TX FIFO 的偏移地址。 (RO)
\label{fielddesc:UARTTXRADDR}\item [UART\_TX\_RADDR] 在 TX FSM 通过 Tx\_FIFO\_Ctrl 读取数据时存储 TX FIFO 的偏移地址。 (RO)
\end{reglist}\end{regdesc}
\end{register}


\begin{register}{H}{UART\_MEM\_RX\_STATUS\_REG}{0x{}0068}\label{regdesc:UARTMEMRXSTATUSREG}
\regfield{(reserved)}{11}{21}{00000000000}%
\regfield{UART\_RX\_WADDR}{10}{11}{{0x{}200}}%
\regfield{(reserved)}{1}{10}{0}%
\regfield{UART\_APB\_RX\_RADDR}{10}{0}{{0x{}200}}%
\reglabel{Reset}\regnewline%
\begin{regdesc}\begin{reglist}
\label{fielddesc:UARTAPBRXRADDR}\item [UART\_APB\_RX\_RADDR] 在软件通过 APB 总线读取 RX FIFO 数据时存储 RX FIFO 的偏移地址。UART0 为 0x{}200，UART1 为 0x{}280，UART2 为 0x{}300。 (RO)
\label{fielddesc:UARTRXWADDR}\item [UART\_RX\_WADDR] 在 Rx\_FIFO\_Ctrl  写 RX FIFO 时存储 RX FIFO 的偏移地址。UART0 为 0x{}200，UART1 为 0x{}280，UART2 为 0x{}300。 (RO)
\end{reglist}\end{regdesc}
\end{register}


\begin{register}{H}{UART\_FSM\_STATUS\_REG}{0x{}006C}\label{regdesc:UARTFSMSTATUSREG}
\regfield{(reserved)}{24}{8}{000000000000000000000000}%
\regfield{UART\_ST\_UTX\_OUT}{4}{4}{{0}}%
\regfield{UART\_ST\_URX\_OUT}{4}{0}{{0}}%
\reglabel{Reset}\regnewline%
\begin{regdesc}\begin{reglist}
\label{fielddesc:UARTSTURXOUT}\item [UART\_ST\_URX\_OUT] 接收器的状态字段。 (RO)
\label{fielddesc:UARTSTUTXOUT}\item [UART\_ST\_UTX\_OUT] 发送器的状态字段。 (RO)
\end{reglist}\end{regdesc}
\end{register}


\begin{register}{H}{UART\_LOWPULSE\_REG}{0x{}0028}\label{regdesc:UARTLOWPULSEREG}
\regfield{(reserved)}{20}{12}{00000000000000000000}%
\regfield{UART\_LOWPULSE\_MIN\_CNT}{12}{0}{{0x{}fff}}%
\reglabel{Reset}\regnewline%
\begin{regdesc}\begin{reglist}
\label{fielddesc:UARTLOWPULSEMINCNT}\item [UART\_LOWPULSE\_MIN\_CNT] 存储低电平脉冲的最短持续时间， 用于波特率检测，单位是 APB\_CLK 时钟周期。 (RO)
\end{reglist}\end{regdesc}
\end{register}


\begin{register}{H}{UART\_HIGHPULSE\_REG}{0x{}002C}\label{regdesc:UARTHIGHPULSEREG}
\regfield{(reserved)}{20}{12}{00000000000000000000}%
\regfield{UART\_HIGHPULSE\_MIN\_CNT}{12}{0}{{0x{}fff}}%
\reglabel{Reset}\regnewline%
\begin{regdesc}\begin{reglist}
\label{fielddesc:UARTHIGHPULSEMINCNT}\item [UART\_HIGHPULSE\_MIN\_CNT] 存储最长高电平脉冲持续时间。 用于波特率检测，单位是 APB\_CLK 时钟周期。 (RO)
\end{reglist}\end{regdesc}
\end{register}


\begin{register}{H}{UART\_RXD\_CNT\_REG}{0x{}0030}\label{regdesc:UARTRXDCNTREG}
\regfield{(reserved)}{22}{10}{0000000000000000000000}%
\regfield{UART\_RXD\_EDGE\_CNT}{10}{0}{{0x{}0}}%
\reglabel{Reset}\regnewline%
\begin{regdesc}\begin{reglist}
\label{fielddesc:UARTRXDEDGECNT}\item [UART\_RXD\_EDGE\_CNT] 存储 RXD 沿变化的次数。 用于波特率检测。 (RO)
\end{reglist}\end{regdesc}
\end{register}


\begin{register}{H}{UART\_POSPULSE\_REG}{0x{}0070}\label{regdesc:UARTPOSPULSEREG}
\regfield{(reserved)}{20}{12}{00000000000000000000}%
\regfield{UART\_POSEDGE\_MIN\_CNT}{12}{0}{{0x{}fff}}%
\reglabel{Reset}\regnewline%
\begin{regdesc}\begin{reglist}
\label{fielddesc:UARTPOSEDGEMINCNT}\item [UART\_POSEDGE\_MIN\_CNT] 存储两个上升沿之间的最小输入时钟计数值。 用于波特率检测。 (RO)
\end{reglist}\end{regdesc}
\end{register}


\begin{register}{H}{UART\_NEGPULSE\_REG}{0x{}0074}\label{regdesc:UARTNEGPULSEREG}
\regfield{(reserved)}{20}{12}{00000000000000000000}%
\regfield{UART\_NEGEDGE\_MIN\_CNT}{12}{0}{{0x{}fff}}%
\reglabel{Reset}\regnewline%
\begin{regdesc}\begin{reglist}
\label{fielddesc:UARTNEGEDGEMINCNT}\item [UART\_NEGEDGE\_MIN\_CNT] 存储两个下降沿之间的最小输入时钟计数值。 用于波特率检测。 (RO)
\end{reglist}\end{regdesc}
\end{register}


\begin{register}{H}{UART\_AT\_CMD\_PRECNT\_REG}{0x{}0050}\label{regdesc:UARTATCMDPRECNTREG}
\regfield{(reserved)}{16}{16}{0000000000000000}%
\regfield{UART\_PRE\_IDLE\_NUM}{16}{0}{{0x{}901}}%
\reglabel{Reset}\regnewline%
\begin{regdesc}\begin{reglist}
\label{fielddesc:UARTPREIDLENUM}\item [UART\_PRE\_IDLE\_NUM] 配置接收器接收第一个 AT\_CMD 字符前的空闲时间，单位是比特时间（即传输一个比特所需的时间）。 (R/W)
\end{reglist}\end{regdesc}
\end{register}


\begin{register}{H}{UART\_AT\_CMD\_POSTCNT\_REG}{0x{}0054}\label{regdesc:UARTATCMDPOSTCNTREG}
\regfield{(reserved)}{16}{16}{0000000000000000}%
\regfield{UART\_POST\_IDLE\_NUM}{16}{0}{{0x{}901}}%
\reglabel{Reset}\regnewline%
\begin{regdesc}\begin{reglist}
\label{fielddesc:UARTPOSTIDLENUM}\item [UART\_POST\_IDLE\_NUM] 配置最后一个 AT\_CMD 字符和后续数据的间隔时间，单位是比特时间（即传输一个比特所需的时间）。 (R/W)
\end{reglist}\end{regdesc}
\end{register}


\begin{register}{H}{UART\_AT\_CMD\_GAPTOUT\_REG}{0x{}0058}\label{regdesc:UARTATCMDGAPTOUTREG}
\regfield{(reserved)}{16}{16}{0000000000000000}%
\regfield{UART\_RX\_GAP\_TOUT}{16}{0}{{11}}%
\reglabel{Reset}\regnewline%
\begin{regdesc}\begin{reglist}
\label{fielddesc:UARTRXGAPTOUT}\item [UART\_RX\_GAP\_TOUT] 配置 AT\_CMD 字符的间隔时间，单位是比特时间（即传输一个比特所需的时间）。 (R/W)
\end{reglist}\end{regdesc}
\end{register}


\begin{register}{H}{UART\_AT\_CMD\_CHAR\_REG}{0x{}005C}\label{regdesc:UARTATCMDCHARREG}
\regfield{(reserved)}{16}{16}{0000000000000000}%
\regfield{UART\_CHAR\_NUM}{8}{8}{{0x{}3}}%
\regfield{UART\_AT\_CMD\_CHAR}{8}{0}{{0x{}2b}}%
\reglabel{Reset}\regnewline%
\begin{regdesc}\begin{reglist}
\label{fielddesc:UARTATCMDCHAR}\item [UART\_AT\_CMD\_CHAR] 配置 AT\_CMD 字符的内容。 (R/W)
\label{fielddesc:UARTCHARNUM}\item [UART\_CHAR\_NUM] 配置接收器接收连续 AT\_CMD 字符的个数。 (R/W)
\end{reglist}\end{regdesc}
\end{register}


\begin{register}{H}{UART\_DATE\_REG}{0x{}007C}\label{regdesc:UARTDATEREG}
\regfield{UART\_DATE}{32}{0}{{0x{}2008270}}%
\reglabel{Reset}\regnewline%
\begin{regdesc}\begin{reglist}
\label{fielddesc:UARTDATE}\item [UART\_DATE] 版本控制寄存器。 (R/W)
\end{reglist}\end{regdesc}
\end{register}


\begin{register}{H}{UART\_ID\_REG}{0x{}0080}\label{regdesc:UARTIDREG}
\regfield{UART\_REG\_UPDATE}{1}{31}{{0}}%
\regfield{UART\_UPDATE\_CTRL}{1}{30}{{1}}%
\regfield{UART\_ID}{30}{0}{{0x{}000500}}%
\reglabel{Reset}\regnewline%
\begin{regdesc}\begin{reglist}
\label{fielddesc:UARTID}\item [UART\_ID] 配置 UART\_ID。 (R/W)
\label{fielddesc:UARTUPDATECTRL}\item [UART\_UPDATE\_CTRL] 用于控制同步模式。0：寄存器配置后，软件需向 UART\_REG\_UPDATE 写 1 同步寄存器。；1：寄存器自动同步至 UART Core 时钟域。 (R/W)
\label{fielddesc:UARTREGUPDATE}\item [UART\_REG\_UPDATE] 软件向该字段写 1，将寄存器值同步到 UART Core 时钟域。该字段在同步完成后由硬件自清。 (R/W/SC)
\end{reglist}\end{regdesc}
\end{register}
